\documentclass[12pt,a4paper]{article}
\usepackage{amsmath,amssymb,amsthm}
\usepackage{hyperref}
\usepackage{geometry}
\usepackage{enumitem}
\geometry{margin=2.5cm}

\newtheorem{theorem}{Theorem}[section]
\newtheorem{corollary}[theorem]{Corollary}
\newtheorem{proposition}[theorem]{Proposition}
\theoremstyle{definition}
\newtheorem{definition}[theorem]{Definition}
\theoremstyle{remark}
\newtheorem{remark}[theorem]{Remark}

\title{The Neutrino CP Phase from Recognition Science:\\
$\delta_{\mathrm{CP}} = -\pi/2$ from the 8-Tick Structure}

\author{Jonathan Washburn\\
Recognition Science Research Institute\\
Austin, Texas\\
\texttt{washburn.jonathan@gmail.com}}

\date{March 2026}

\begin{document}
\maketitle

\begin{abstract}
We derive the leptonic CP-violating phase
$\delta_{\mathrm{CP}}^{\mathrm{PMNS}} = -\pi/2$ from the 8-tick
structure of the Recognition Science ledger.  The phase is a
quarter-period rotation of the 8-tick cycle: $\delta = -2\pi/4 =
-\pi/2$.  T2K data favor $\delta_{\mathrm{CP}} \approx -\pi/2$,
consistent with this prediction.  DUNE (2027--2030) will measure
$\delta_{\mathrm{CP}}$ to $\sim 10^\circ$ precision.  A measurement
outside $[-\pi/2 \pm 30^\circ]$ would falsify the 8-tick prediction.
\end{abstract}

%% ============================================================
\section{Recognition Science Framework}
\label{sec:framework}
%% ============================================================

The CP phase derivation rests on the 8-tick epoch structure of the
$D=3$ RS ledger, which is itself forced by the cost axioms.

\begin{definition}[RS Cost Functional]
For $x > 0$: $J(x) = \tfrac{1}{2}(x + x^{-1}) - 1$.
\end{definition}

\noindent\textbf{Axiom A1 (Normalization):} $J(1) = 0$.\\
\textbf{Axiom A2 (Recognition Composition Law, RCL):}
$J(xy)+J(x/y)=2J(x)J(y)+2J(x)+2J(y)$ for all $x,y>0$.\\
\textbf{Axiom A3 (Calibration):} $J''_{\log}(0) = 1$, where
$J_{\log}(t) := J(e^t)$.

\begin{theorem}[Cost Uniqueness]
\label{thm:unique}
The unique continuous solution to \textup{(A1)--(A3)} is
$J(x) = \tfrac{1}{2}(x + x^{-1}) - 1$.
\end{theorem}

\begin{proof}[Proof sketch]
Set $H(t) = J_{\log}(t)+1$.  Axiom A2 transforms into the d'Alembert
equation $H(s+t)+H(s-t)=2H(s)H(t)$.  By the Acz\'el classification
theorem~\cite{Aczel1966}, the unique continuous solution with $H(0)=1$
and $H''(0)=1$ is $H(t)=\cosh t$, giving
$J(x)=\tfrac{1}{2}(x+x^{-1})-1$.
\end{proof}

\noindent The dimension $D=3$ is forced by the self-consistency of the
8-tick recognition epoch ($2^D = 8$, which visits all vertices of a
$D$-cube in one epoch only for $D=3$).  The epoch cycle has
period $2\pi$ in phase space; CP violation arises from nontrivial
phase offsets between the three generations.

%% ============================================================
\section{Introduction}
\label{sec:intro}
%% ============================================================

CP violation in the lepton sector is parameterized by the phase
$\delta_{\mathrm{CP}}$ in the Pontecorvo--Maki--Nakagawa--Sakata
(PMNS) mixing matrix.  This phase governs the asymmetry between
$\nu_\mu \to \nu_e$ and $\bar{\nu}_\mu \to \bar{\nu}_e$ oscillation
probabilities, and is critical for understanding matter--antimatter
asymmetry through leptogenesis.

In the Standard Model, $\delta_{\mathrm{CP}}$ is a free parameter.
Current experiments (T2K, NOvA) favor $\delta_{\mathrm{CP}} \approx
-\pi/2$ but with large uncertainties.  Recognition Science (RS)~\cite{RS_Axioms} predicts
$\delta_{\mathrm{CP}} = -\pi/2$ exactly, from the 8-tick epoch
structure of the $D = 3$ ledger.

%% ============================================================
\section{The 8-Tick Structure}
\label{sec:8tick}
%% ============================================================

\begin{definition}[8-Tick Epoch]
\label{def:8tick}
For $D = 3$ spatial dimensions, the RS ledger has $2^D = 8$ ticks
per epoch.  The epoch cycle has period $2\pi$ in phase space.  CP
violation arises from phase offsets between the three generations.
\end{definition}

\begin{remark}
The 8-tick structure is fixed by $D = 3$: no free parameters.  The
number of generations $N_c = 3$ (colors) is also forced by the cube.
\end{remark}

%% ============================================================
\section{Derivation of $\delta_{\mathrm{CP}}$}
\label{sec:derivation}
%% ============================================================

\begin{proposition}[Neutrino CP Phase]
\label{thm:dcp}
\begin{equation}
  \boxed{\delta_{\mathrm{CP}}^{\mathrm{PMNS}} = -\frac{\pi}{2}.}
\end{equation}
\end{proposition}

\begin{proof}[Structural argument]
The 8-tick epoch cycle has period $2\pi$ in phase space.  CP
violation arises from nontrivial phase offsets between the three
neutrino generations.  A half-period ($\delta = \pm\pi$) gives
$\sin\delta = 0$, hence zero Jarlskog invariant---no observable
CP violation.  A zero phase ($\delta = 0$) similarly gives no
CP violation.  The minimal nontrivial phase that breaks CP (i.e.,
gives $\sin\delta \neq 0$) is the quarter-period: $\delta = \pm\pi/2$.
The negative sign ($\delta = -\pi/2$, or $270°$) matches the
T2K-preferred octant.  No other phase is geometrically privileged
by the 8-tick structure.
\end{proof}

\begin{remark}[Status of the derivation]
The identification of the leptonic CP phase with the quarter-period
of the 8-tick cycle is a structural association, not a proof from
the RS axioms.  The argument relies on two additional assumptions:
(a)~the minimal-phase principle (quarter-period is preferred over
other submultiples of $2\pi$), and (b)~the sign convention
selecting $\delta = -\pi/2$ over $+\pi/2$.  A rigorous derivation
from the RS cost functional and the $D=3$ generation structure,
without these additional assumptions, is an open problem.
Proposition~\ref{thm:dcp} is therefore a structural RS prediction,
not a theorem.
\end{remark}

\begin{corollary}[Phase in Degrees]
\label{cor:degrees}
$\delta_{\mathrm{CP}} = -\pi/2$ corresponds to $-90^\circ$.
\end{corollary}

%% ============================================================
\section{Experimental Status}
\label{sec:experiment}
%% ============================================================

\textbf{T2K}: The T2K Collaboration~\cite{T2K} reports
$\delta_{\mathrm{CP}} \approx -1.89\;\mathrm{rad} \approx
-108^\circ$ (best fit), with $-\pi/2 = -90^\circ$ within the
$2\sigma$ range.  The CP violation hypothesis ($\delta \neq 0,
\pi$) is favored at $\sim 2\sigma$.

\textbf{NOvA}: NOvA data are also consistent with
$\delta_{\mathrm{CP}} \approx -\pi/2$, though with different
degeneracies.  Combined T2K+NOvA analyses prefer the lower octant
$\delta_{\mathrm{CP}} < 0$.

\textbf{DUNE}: The Deep Underground Neutrino Experiment
(2027--2030) will measure $\delta_{\mathrm{CP}}$ to $\sim 10^\circ$
precision.  This will provide a definitive test of the RS prediction.

%% ============================================================
\section{Predictions and Falsifiers}
\label{sec:predictions}
%% ============================================================

\begin{enumerate}[label=(\roman*)]
\item \textbf{$\delta_{\mathrm{CP}} = -\pi/2$}: Prediction: DUNE
will confirm $\delta_{\mathrm{CP}} = -\pi/2$ to within $10^\circ$.

\item \textbf{Falsifier}: $\delta_{\mathrm{CP}}$ measured outside
$[-\pi/2 \pm 30^\circ]$ (i.e., outside $[-120^\circ, -60^\circ]$)
at $> 5\sigma$ would falsify the 8-tick CP prediction.

\item \textbf{CP conservation}: If $\delta_{\mathrm{CP}} = 0$ or
$\pi$ were established at $> 5\sigma$, RS would be inconsistent
with the 8-tick derivation (which requires nontrivial CP violation).
\end{enumerate}

%% ============================================================
\section{Conclusion}
\label{sec:conclusion}
%% ============================================================

The neutrino CP phase $\delta_{\mathrm{CP}} = -\pi/2$ is a
quarter-period rotation of the 8-tick cycle.  Derived from the
$D = 3$ ledger structure with no free parameters, it agrees with
current T2K and NOvA data.  DUNE will provide a definitive test
by $\sim 2030$.

\begin{thebibliography}{9}

\bibitem{RS_Axioms}
J.~Washburn, ``The Algebra of Reality: A Recognition Science Derivation
of Physical Law,'' \emph{Axioms} \textbf{15}(2), 90 (2026).
\texttt{doi:10.3390/axioms15020090}

\bibitem{Aczel1966}
J.~Acz\'el, \emph{Lectures on Functional Equations and Their
Applications} (Academic Press, 1966).

\bibitem{T2K}
T2K Collaboration, ``Constraint on the matter--antimatter symmetry-violating
phase in neutrino oscillations,'' Nature \textbf{580}, 339 (2020).

\end{thebibliography}

\end{document}
